\SourceLeaf{275}%
independent geometry, which he had pursued on one occasion and called anti-Euclidean geometry.

According to the deepest conviction he often expressed, \textsc{Gauss} regarded the three dimensions of space as a particular property of the human mind. In a humorous mood he once called people who could not see this Boeotians. We can imagine ourselves, he said, in the position of beings conscious of only two dimensions. Beings of a higher order might look down on us in a similar way. He went on in jest to say that he had set aside certain problems here that he intended to treat geometrically later, in a higher state.\}

\Needspace{6\baselineskip}
\begin{center}\textsc{Remark.}\end{center}

The paper mentioned by \textsc{Sartorius von Waltershausen} was found among \textsc{Gauss}’s papers. It deals essentially only with explaining the concept of number and the four rules of arithmetic.

\hfill \textsc{Stäckel}.
